% =====================================================================
%  Sampling method and detection rules
%  Drop-in section for paper_draft.tex:  % =====================================================================
%  Sampling method and detection rules
%  Drop-in section for paper_draft.tex:  % =====================================================================
%  Sampling method and detection rules
%  Drop-in section for paper_draft.tex:  % =====================================================================
%  Sampling method and detection rules
%  Drop-in section for paper_draft.tex:  \input{method_sampling_rules}
%
%  Assumes the preamble already loads: booktabs, amsmath, amssymb.
%  Uses no algorithm package: the pseudocode sits in a plain figure float.
%  Cites: klick2016, bozkurt:dissecting:2018, ripestat  (all in references.bib)
%
%  Every figure below is produced by the pipeline in
%  experiments/15_national_block_routing_census/.
% =====================================================================

\section{Measurement Method}
\label{sec:method}

\subsection{Universe}
\label{sec:universe}

``Pakistan's address space'' has two definitions that disagree, and neither contains the
other. The \emph{registry view} is the space recorded against Pakistan at the RIR. The
\emph{routing view} is the space a Pakistani AS originates in BGP. Registration is
paperwork; origination is a deliberate operational act, and BGP does not require the two
to agree.

We take the union. Of the resulting space, \textbf{252{,}672 addresses (4.4\%) are
originated by a Pakistani AS but registered elsewhere}, typically leased space announced
as a more-specific inside a covering announcement the registrant still originates. In the
other direction, \textbf{1{,}209 /24 blocks} are registered to Pakistan and originated by
nobody. A study defining its universe by registration alone discards the first set; one
defining it by origination alone discards the second.

\begin{table}[t]
\centering
\caption{Measurement universe.}
\label{tab:universe}
\begin{tabular}{lr}
\toprule
Addresses (registry $\cup$ routing) & 5{,}774{,}336 \\
Announced or registered prefixes    & 932 \\
/24-equivalent blocks               & 22{,}556 \\
Networks with announced space       & 352 \\
\bottomrule
\end{tabular}
\end{table}

Prefixes are subdivided into /24-equivalent \emph{blocks} so that sampling density is
uniform: without this a /16 and a /24 would receive equal attention despite differing by a
factor of 256. The mapping is one-to-many, not one-to-one: 686 of the 932 prefixes are
larger than a /24. A per-block count is therefore not a per-prefix count.

Origination is queried through RIPEstat~\cite{ripestat}, which is fed by RIS. We note that
origination is \emph{not} a boolean: a prefix is originated \emph{as seen by $n$ of the
324 IPv4 RIS peers}, and RIPEstat endpoints that report a single visibility-filtered
origin can disagree with the full origin history for the same prefix at the same instant.
Attribution is therefore taken from origin history rather than from a filtered view.

\subsection{Adaptive liveness sampling}
\label{sec:sampling}

Exhaustive probing of 5.77M addresses is neither necessary nor courteous. We follow the
prefix-as-unit and density formulation of TASS~\cite{klick2016}, but \emph{invert its
rule}: TASS discards low-density prefixes to shrink a full-space scan, whereas we retain
every prefix and reduce effort within each one.

A \emph{check} tests one address by trying, in order, ICMP echo, TCP connect to port~80,
and TCP connect to port~443, stopping at the first response. A refused TCP connection
counts as alive, since something was present to refuse it. One check may thus emit up to
three packets. A \emph{draw} is a batch of $D=8$ addresses.

\begin{figure}[t]
\small
\hrule\vspace{4pt}
\begin{tabbing}
xx\=xx\=xx\=xx\kill
\textbf{for each} block $B$: \\
\> $S \leftarrow$ addresses of $B$ already probed by any prior scan \\
\> $L \leftarrow$ live hosts already known in $B$ \\
\> $P \leftarrow \mathrm{shuffle}\!\left(\mathrm{hosts}(B) \setminus S\right)$
   \quad \textit{(without replacement)} \\
\> $r \leftarrow 0$ \\
\> \textbf{while} $P \neq \emptyset$ \textbf{and} $|S| < M$ \textbf{and} $L < K$: \\
\> \> $b \leftarrow$ next $D$ addresses of $P$; \quad $r \leftarrow r+1$ \\
\> \> \textbf{for} $a \in b$: \\
\> \> \> $S \leftarrow S \cup \{a\}$;\quad
        \textbf{if} $\mathrm{check}(a)$ \textbf{then} $L \leftarrow L+1$ \\
\> \> \> record $(a, \mathrm{result})$ \quad \textit{(negative results retained)} \\
\> \> \textbf{if} $L = 0$ \textbf{and} $r \geq 2$ \textbf{then break} \\
\> emit $L$ for $B$
\end{tabbing}
\vspace{2pt}\hrule
\caption{Adaptive per-block sampling. $D=8$, $K=8$, $M=64$.}
\label{alg:sample}
\end{figure}

Sampling is \textbf{without replacement}: $S$ accumulates across draws and across scans,
so no address is ever probed twice. Draws accumulate rather than replace, and negative
results are retained because the miss-rate estimates of \S\ref{sec:limits} depend on them.
The target is re-tested per draw rather than per address, so a block may terminate above
$K$; 17.2\% of live blocks did.

\paragraph{Choice of $K$.} $K=8$ is measured rather than assumed, and the criterion is
\emph{split-block detection}, not redundancy. Over 1{,}758 (block, vantage) pairs that
received a full $K=8$ in a prior round, \textbf{20.0\% of blocks were split}: some probed
addresses detoured abroad while others remained domestic, so ``the route to a block'' is
not always well defined. On the 79.4\% of pairs that were unambiguous, a randomly chosen
$K=2$ reproduces the $K=8$ verdict 98.6\% of the time. A small $K$ therefore recovers the
\emph{verdict} while destroying the \emph{picture}: with two samples there is nothing to
compare against, and a split block is silently recorded as whichever pair was drawn. A
20\% split rate is too large to sample away.

\paragraph{Cost of the stopping rule.} The two-empty-draw rule is the design's principal
source of missed occupancy, and its cost is calculable. Sampling without replacement from
$N=254$ usable addresses, the probability that all $d=16$ draws miss a block containing
$\ell$ live hosts is hypergeometric:
\begin{equation}
\Pr[\text{all } d \text{ dead}] \;=\; \prod_{i=0}^{d-1} \frac{N-\ell-i}{N-i}.
\label{eq:hyper}
\end{equation}

\begin{table}[t]
\centering
\caption{Probability that the two-empty-draw rule abandons a populated block,
from Eq.~\eqref{eq:hyper}.}
\label{tab:hyper}
\begin{tabular}{rrr}
\toprule
Live hosts $\ell$ & Density & $\Pr[\text{abandoned}]$ \\
\midrule
1  & 0.4\%  & 93.7\% \\
2  & 0.8\%  & 87.8\% \\
5  & 2.0\%  & 72.0\% \\
8  & 3.1\%  & 59.0\% \\
13 & 5.1\%  & 42.0\% \\
26 & 10.2\% & 16.8\% \\
51 & 20.1\% & \phantom{0}2.4\% \\
\bottomrule
\end{tabular}
\end{table}

A block holding eight live hosts is abandoned 59\% of the time. \textbf{Block-level
coverage figures are therefore floors, and substantially so}: blocks we record as empty
are not reliably empty. This does not bias the routing results of \S\ref{sec:rules},
which condition on blocks that were successfully traced, but it forbids any claim about
the extent of unused address space.

\begin{table}[t]
\centering
\caption{Where sampling effort went. Effort follows signal: a populated block
receives roughly four times the attention of an empty one.}
\label{tab:effort}
\begin{tabular}{lrrr}
\toprule
 & Blocks & Checks & Mean/block \\
\midrule
Recorded empty & 16{,}669 & 364{,}972 & 21.9 \\
Contained life & \phantom{0}5{,}887 & 479{,}130 & 81.4 \\
\midrule
Total          & 22{,}556 & 844{,}102 & 37.4 \\
\bottomrule
\end{tabular}
\end{table}

\textbf{844{,}102 addresses were checked, 14.6\% of the universe, yielding 43{,}737 live
hosts.} The answer rate is 5.2\% \emph{of what was tested}; the remaining 4{,}930{,}234
addresses were never probed, and the count must not be read against the full space.

\subsection{Route collection and trace selection}
\label{sec:routes}

Every live host was traced with ICMP, incrementing the hop limit and terminating after
five consecutive silent hops. Of \textbf{43{,}765} traces attempted, \textbf{36{,}214
(82.7\%)} reached their target.

A trace is retained only if it \emph{reached} its target, so that the endpoint is known,
and at least five hops responded, so that the path has discernible shape. At most eight
traces are kept per block, ranked by fewest timeouts. This yields \textbf{34{,}191 traces
across 5{,}508 blocks in 217 networks}, comprising \textbf{217{,}255 responding hops},
each resolved to a country and an operator.

The five-hop gate is a judgement rather than a derived quantity, and it does not fall
evenly: it binds almost entirely on one large network whose paths are of ordinary length
but answer at half the usual rate (median 5 responding hops against 8 elsewhere, 48\%
silent slots against 31\%).

\section{Detection Rules}
\label{sec:rules}

A rule here is \textbf{a procedure, not a threshold}. Each is stated as a claim that could
be false, tested against a vantage point's own data, and applied only where it survives. A
rule proven at one vantage says nothing about another until retested there. Consequently
every numeric bound below is an \emph{output} of the procedure at our vantage, not a
constant of the method.

\paragraph{The single constant.} The speed of light in fibre, $c_f = 204{,}000$\,km/s,
taken from Bozkurt et al.~\cite{bozkurt:dissecting:2018}. It is a physical property, and
it appears in exactly one rule. A detector cannot derive physics from its own traffic.

We emphasise what is \emph{not} used. Bozkurt et al. also give a widely quoted rule of
thumb, that internet latency is approximately line-of-sight distance $\times\,2.1$ divided
by $c_f$. That is a \emph{typical-case estimator} and not an upper bound; the same work
reports that only 11\% of real fibre links fall within 25\% of the latency their length
predicts, and that co-located server pairs frequently differ by 10--30\,ms. Using it as a
domestic ceiling is unsound. Worse, at national scale it is the wrong functional form:
regressing clean minimum RTT on straight-line distance over 346 domestic destinations
yields $R^2 = 0.013$ with a slightly negative slope. With a median domestic path of
$\approx$470\,km (4.6\,ms of propagation) against 25--35\,ms of fixed access overhead,
distance does not drive domestic latency, and no multiplier over distance can describe it.

\subsection{R0: the measurement can support a latency rule}

\textbf{Claim.} The RTT data is a measurement rather than an artefact of the measuring
process. \textbf{Test.} Two checks requiring no external reference: does RTT ever
\emph{decrease} with hop distance, which propagation cannot produce; and does a fixed
router return a stable value under repetition? \textbf{Gate.} If either fails, no latency
rule may be evaluated on that data; the measurement is repaired rather than the thresholds
retuned.

Our initial route sweep failed both, having emitted one packet per hop limit under 120-way
concurrency: 22\% of adjacent hop pairs showed RTT falling by more than 20\,ms with
increasing distance, and a single fixed router spanned 4--853\,ms. Re-measurement at 6--8
threads with 20--30 packets per address yields 1--2\,ms of interquartile dispersion and
values a median of 30--41\,ms lower. All latency figures below derive from the
re-measurement. \textbf{Routers are measured in-path}, with packets addressed to the
original destination under a constrained hop limit, never by addressing the router
directly: routers deprioritise traffic addressed to themselves, and the route \emph{to} a
router need not coincide with the route \emph{through} it.

\subsection{R1: the vantage's access chain}

\textbf{Claim.} All traces from vantage $v$ share an initial router sequence $C(v)$ until
routes diverge. \textbf{Test.} The dominant address at each early hop position.
Parameter-free, since the divergence point is measured.

At our vantage, three routers appear in $\geq$98.8\% of traces, diverging at hop~6, and
cost $B_{\mathrm{access}} = 20$\,ms. This is our own connection and is subtracted before
any latency is attributed to a destination.

\subsection{R2: an address is not where it is registered}

\textbf{Claim.} If address $X$ registered in country $C$ were physically in $C$, no
Pakistani vantage could observe it faster than $2\,d(v,C)/c_f$. \textbf{Test.} Clean
in-path minimum RTT against that floor, with $d$ computed from per-address geolocation and
the vantage's own geolocated position. No routing factor is applied, so a falsification is
a physical impossibility rather than a judgement. \textbf{Gate.} Foreign claims only;
domestic distances are too short for physics to discriminate.

The minimum is the correct estimator here because queuing, rate limiting and slow ICMP
generation are strictly additive, so the minimum of a controlled burst is the closest
available approach to propagation cost.

Of the foreign-registered addresses observed as hops, \textbf{79 are falsified}, 341 are
upheld, and 78 are undecidable because their country's floor lies inside the domestic
range. The falsified set includes US Department of Defense addresses responding in
$\approx$3\,ms, together with Cogent, Cloudflare, AT\&T and T-Mobile ranges: Pakistani
routers numbered from foreign address space. Absent this step these hops would be counted
as departures and would fabricate thousands of spurious detours.

\textbf{R2 falsifies a claim; it does not establish a location.} ``Not in Virginia'' is
not ``in Pakistan'', which is what R4$'$ decides.

\subsection{R3: the vantage can observe paths, under this protocol}

\textbf{Claim.} Traces from $v$ contain public addresses. \textbf{Gate.} Established per
protocol, never universally: a probe is usable \emph{under a protocol}, and in earlier
data the same device yielded no public hop in 100\% of TCP traces and 11\% of ICMP traces.
Here, 3 of 43{,}765 traces contained no public router.

\subsection{R4$'$: the vantage's domestic normal is its own}

\textbf{Claim.} A destination whose clean RTT lies far outside this vantage's own domestic
distribution did not traverse a domestic path.

\textbf{Test.} Construct the reference set from destinations whose observed path is
entirely domestic; measure each cleanly and retain the minimum; flag any value beyond the
Tukey far fence,
\begin{equation}
\tau \;=\; Q_3 + 3\,(Q_3 - Q_1).
\label{eq:fence}
\end{equation}

Equation~\eqref{eq:fence} is used because it is a standard, scale-free outlier definition
that no one selects to fit a dataset. A slow vantage obtains a high baseline \emph{and} a
high fence, so \textbf{the procedure transfers between vantages without retuning}, which a
millisecond threshold cannot.

\textbf{Gates.} At least $\approx$100 reference destinations, or the distribution's shape
is noise; and if more than half the reference set is itself detouring, the baseline
\emph{is} the detour and the rule is rejected for that vantage.

\textbf{Scope.} The reference set is selected topologically, so R4$'$ can only recover
detours the topological test missed. It is a second filter beneath the first, not a
replacement.

\begin{table}[t]
\centering
\caption{R4$'$ evaluated at our vantage. The bound is an output of the procedure, not a
constant of the method.}
\label{tab:r4}
\begin{tabular}{lr}
\toprule
Reference destinations              & 347 \\
Domestic median                     & 27\,ms \\
$Q_1$ / $Q_3$                       & 24 / 38\,ms \\
\textbf{Fence $\tau$}               & \textbf{80\,ms} \\
Domestic above $\tau$               & 1.4\% \\
Known detours below $\tau$          & 0 of 10 \\
\bottomrule
\end{tabular}
\end{table}

The two distributions separate empirically: the highest domestic value below the fence is
69\,ms and the lowest tromboning destination is 90\,ms.

\subsection{R5: the trace left the country and returned}

Evaluated only after R0--R4$'$ are settled for that vantage under that protocol. Anything
attributable to an earlier rule is subtracted first, and a detour is the residue:

\begin{enumerate}\itemsep2pt
\item subtract $C(v)$; the access chain cannot be a detour;
\item apply R2; a foreign registry is not evidence, an impossible RTT is;
\item apply R4$'$ to whatever R2 falsified, since ``not in $C$'' is not ``domestic'';
\item any surviving foreign hop constitutes a departure.
\end{enumerate}

A hop on a \emph{foreign internet exchange fabric} counts as a departure: the address
belongs to the exchange rather than to either peer, but the packet is physically abroad.

\textbf{Result: 930 of 34{,}191 traces (2.72\%), across 156 blocks in 49 networks.} This
is a \textbf{floor}. With the latency arm confined to R4$'$'s scope, a detour is visible
only when a foreign router responds; detours through silent routers or inside MPLS tunnels
are recorded as domestic.

\paragraph{Blocks, not traces, are the unit of evidence.} Tromboning is a property of the
destination prefix: 146 of 156 affected blocks are all-or-nothing, every selected host
detouring or none, and 143 of 156 exit via a single common address. Eight traces into one
block are eight observations of one routing decision. Confidence intervals must therefore
be computed over blocks and clustered by network.

\subsection{R6: the verdict predicts what it was not given}

\textbf{Claim.} If R5 detects a real phenomenon, its verdicts should separate on a
quantity the topological test never observed.

\textbf{Test.} Compare clean minimum RTT between destinations R5 labels domestic and those
it labels tromboning. Destinations labelled domestic have a median of \textbf{27\,ms};
those labelled tromboning, \textbf{143\,ms}. The topological classifier had no access to
latency, so the 116\,ms separation is independent corroboration rather than a restatement.

\textbf{Limit.} The tromboning arm of this comparison holds 10 destinations. It is
corroboration, not a significance test, and warrants repetition on a larger deliberate
sample.

\section{Limitations}
\label{sec:limits}

\textbf{Counts are floors with unequal tightness.} Adaptive sampling caps at 64 addresses
per block, so no block reports more than 64 regardless of occupancy. Independently,
re-probing recorded-dead addresses returns $\approx$35\% alive on one large network
against 22--25\% elsewhere. Densities are therefore comparable within a network over time,
not between networks, and no single national correction is admissible.

\textbf{Liveness is vantage-relative.} In a controlled retest, $\approx$9\% of addresses
confirmed alive from one Pakistani network did not respond from another, at every
concurrency tested. Reported counts are a union over two vantages covering
\emph{disjoint} address sets, not a two-vantage measurement of the same addresses.

\textbf{Route collection is effectively single-vantage.} 43{,}671 of 43{,}765 traces
egress one network. Reach rates are properties of that vantage's view, not of the
destination network alone.

\textbf{Proximate foreign destinations are undetectable by latency.} Domestic paths run at
a 27\,ms median while the physical floors to the nearest Gulf hubs are 4--7\,ms. A Gulf
detour and a slow domestic path are indistinguishable by RTT at any threshold; only
topology can recover them.

\textbf{Selection favours dense blocks.} A block requires live hosts before it can be
traced, and the retention gate then prefers responsive paths. A control comparison against
sparsely populated blocks has not been run, so the magnitude of this bias is unknown
rather than small.

\textbf{Remaining limits.} Traceroute observes the forward path only; return paths may
differ and are invisible. MPLS tunnels conceal interior hops. The study is IPv4 only,
covers a single country over three days, and no operator has confirmed any result.

%
%  Assumes the preamble already loads: booktabs, amsmath, amssymb.
%  Uses no algorithm package: the pseudocode sits in a plain figure float.
%  Cites: klick2016, bozkurt:dissecting:2018, ripestat  (all in references.bib)
%
%  Every figure below is produced by the pipeline in
%  experiments/15_national_block_routing_census/.
% =====================================================================

\section{Measurement Method}
\label{sec:method}

\subsection{Universe}
\label{sec:universe}

``Pakistan's address space'' has two definitions that disagree, and neither contains the
other. The \emph{registry view} is the space recorded against Pakistan at the RIR. The
\emph{routing view} is the space a Pakistani AS originates in BGP. Registration is
paperwork; origination is a deliberate operational act, and BGP does not require the two
to agree.

We take the union. Of the resulting space, \textbf{252{,}672 addresses (4.4\%) are
originated by a Pakistani AS but registered elsewhere}, typically leased space announced
as a more-specific inside a covering announcement the registrant still originates. In the
other direction, \textbf{1{,}209 /24 blocks} are registered to Pakistan and originated by
nobody. A study defining its universe by registration alone discards the first set; one
defining it by origination alone discards the second.

\begin{table}[t]
\centering
\caption{Measurement universe.}
\label{tab:universe}
\begin{tabular}{lr}
\toprule
Addresses (registry $\cup$ routing) & 5{,}774{,}336 \\
Announced or registered prefixes    & 932 \\
/24-equivalent blocks               & 22{,}556 \\
Networks with announced space       & 352 \\
\bottomrule
\end{tabular}
\end{table}

Prefixes are subdivided into /24-equivalent \emph{blocks} so that sampling density is
uniform: without this a /16 and a /24 would receive equal attention despite differing by a
factor of 256. The mapping is one-to-many, not one-to-one: 686 of the 932 prefixes are
larger than a /24. A per-block count is therefore not a per-prefix count.

Origination is queried through RIPEstat~\cite{ripestat}, which is fed by RIS. We note that
origination is \emph{not} a boolean: a prefix is originated \emph{as seen by $n$ of the
324 IPv4 RIS peers}, and RIPEstat endpoints that report a single visibility-filtered
origin can disagree with the full origin history for the same prefix at the same instant.
Attribution is therefore taken from origin history rather than from a filtered view.

\subsection{Adaptive liveness sampling}
\label{sec:sampling}

Exhaustive probing of 5.77M addresses is neither necessary nor courteous. We follow the
prefix-as-unit and density formulation of TASS~\cite{klick2016}, but \emph{invert its
rule}: TASS discards low-density prefixes to shrink a full-space scan, whereas we retain
every prefix and reduce effort within each one.

A \emph{check} tests one address by trying, in order, ICMP echo, TCP connect to port~80,
and TCP connect to port~443, stopping at the first response. A refused TCP connection
counts as alive, since something was present to refuse it. One check may thus emit up to
three packets. A \emph{draw} is a batch of $D=8$ addresses.

\begin{figure}[t]
\small
\hrule\vspace{4pt}
\begin{tabbing}
xx\=xx\=xx\=xx\kill
\textbf{for each} block $B$: \\
\> $S \leftarrow$ addresses of $B$ already probed by any prior scan \\
\> $L \leftarrow$ live hosts already known in $B$ \\
\> $P \leftarrow \mathrm{shuffle}\!\left(\mathrm{hosts}(B) \setminus S\right)$
   \quad \textit{(without replacement)} \\
\> $r \leftarrow 0$ \\
\> \textbf{while} $P \neq \emptyset$ \textbf{and} $|S| < M$ \textbf{and} $L < K$: \\
\> \> $b \leftarrow$ next $D$ addresses of $P$; \quad $r \leftarrow r+1$ \\
\> \> \textbf{for} $a \in b$: \\
\> \> \> $S \leftarrow S \cup \{a\}$;\quad
        \textbf{if} $\mathrm{check}(a)$ \textbf{then} $L \leftarrow L+1$ \\
\> \> \> record $(a, \mathrm{result})$ \quad \textit{(negative results retained)} \\
\> \> \textbf{if} $L = 0$ \textbf{and} $r \geq 2$ \textbf{then break} \\
\> emit $L$ for $B$
\end{tabbing}
\vspace{2pt}\hrule
\caption{Adaptive per-block sampling. $D=8$, $K=8$, $M=64$.}
\label{alg:sample}
\end{figure}

Sampling is \textbf{without replacement}: $S$ accumulates across draws and across scans,
so no address is ever probed twice. Draws accumulate rather than replace, and negative
results are retained because the miss-rate estimates of \S\ref{sec:limits} depend on them.
The target is re-tested per draw rather than per address, so a block may terminate above
$K$; 17.2\% of live blocks did.

\paragraph{Choice of $K$.} $K=8$ is measured rather than assumed, and the criterion is
\emph{split-block detection}, not redundancy. Over 1{,}758 (block, vantage) pairs that
received a full $K=8$ in a prior round, \textbf{20.0\% of blocks were split}: some probed
addresses detoured abroad while others remained domestic, so ``the route to a block'' is
not always well defined. On the 79.4\% of pairs that were unambiguous, a randomly chosen
$K=2$ reproduces the $K=8$ verdict 98.6\% of the time. A small $K$ therefore recovers the
\emph{verdict} while destroying the \emph{picture}: with two samples there is nothing to
compare against, and a split block is silently recorded as whichever pair was drawn. A
20\% split rate is too large to sample away.

\paragraph{Cost of the stopping rule.} The two-empty-draw rule is the design's principal
source of missed occupancy, and its cost is calculable. Sampling without replacement from
$N=254$ usable addresses, the probability that all $d=16$ draws miss a block containing
$\ell$ live hosts is hypergeometric:
\begin{equation}
\Pr[\text{all } d \text{ dead}] \;=\; \prod_{i=0}^{d-1} \frac{N-\ell-i}{N-i}.
\label{eq:hyper}
\end{equation}

\begin{table}[t]
\centering
\caption{Probability that the two-empty-draw rule abandons a populated block,
from Eq.~\eqref{eq:hyper}.}
\label{tab:hyper}
\begin{tabular}{rrr}
\toprule
Live hosts $\ell$ & Density & $\Pr[\text{abandoned}]$ \\
\midrule
1  & 0.4\%  & 93.7\% \\
2  & 0.8\%  & 87.8\% \\
5  & 2.0\%  & 72.0\% \\
8  & 3.1\%  & 59.0\% \\
13 & 5.1\%  & 42.0\% \\
26 & 10.2\% & 16.8\% \\
51 & 20.1\% & \phantom{0}2.4\% \\
\bottomrule
\end{tabular}
\end{table}

A block holding eight live hosts is abandoned 59\% of the time. \textbf{Block-level
coverage figures are therefore floors, and substantially so}: blocks we record as empty
are not reliably empty. This does not bias the routing results of \S\ref{sec:rules},
which condition on blocks that were successfully traced, but it forbids any claim about
the extent of unused address space.

\begin{table}[t]
\centering
\caption{Where sampling effort went. Effort follows signal: a populated block
receives roughly four times the attention of an empty one.}
\label{tab:effort}
\begin{tabular}{lrrr}
\toprule
 & Blocks & Checks & Mean/block \\
\midrule
Recorded empty & 16{,}669 & 364{,}972 & 21.9 \\
Contained life & \phantom{0}5{,}887 & 479{,}130 & 81.4 \\
\midrule
Total          & 22{,}556 & 844{,}102 & 37.4 \\
\bottomrule
\end{tabular}
\end{table}

\textbf{844{,}102 addresses were checked, 14.6\% of the universe, yielding 43{,}737 live
hosts.} The answer rate is 5.2\% \emph{of what was tested}; the remaining 4{,}930{,}234
addresses were never probed, and the count must not be read against the full space.

\subsection{Route collection and trace selection}
\label{sec:routes}

Every live host was traced with ICMP, incrementing the hop limit and terminating after
five consecutive silent hops. Of \textbf{43{,}765} traces attempted, \textbf{36{,}214
(82.7\%)} reached their target.

A trace is retained only if it \emph{reached} its target, so that the endpoint is known,
and at least five hops responded, so that the path has discernible shape. At most eight
traces are kept per block, ranked by fewest timeouts. This yields \textbf{34{,}191 traces
across 5{,}508 blocks in 217 networks}, comprising \textbf{217{,}255 responding hops},
each resolved to a country and an operator.

The five-hop gate is a judgement rather than a derived quantity, and it does not fall
evenly: it binds almost entirely on one large network whose paths are of ordinary length
but answer at half the usual rate (median 5 responding hops against 8 elsewhere, 48\%
silent slots against 31\%).

\section{Detection Rules}
\label{sec:rules}

A rule here is \textbf{a procedure, not a threshold}. Each is stated as a claim that could
be false, tested against a vantage point's own data, and applied only where it survives. A
rule proven at one vantage says nothing about another until retested there. Consequently
every numeric bound below is an \emph{output} of the procedure at our vantage, not a
constant of the method.

\paragraph{The single constant.} The speed of light in fibre, $c_f = 204{,}000$\,km/s,
taken from Bozkurt et al.~\cite{bozkurt:dissecting:2018}. It is a physical property, and
it appears in exactly one rule. A detector cannot derive physics from its own traffic.

We emphasise what is \emph{not} used. Bozkurt et al. also give a widely quoted rule of
thumb, that internet latency is approximately line-of-sight distance $\times\,2.1$ divided
by $c_f$. That is a \emph{typical-case estimator} and not an upper bound; the same work
reports that only 11\% of real fibre links fall within 25\% of the latency their length
predicts, and that co-located server pairs frequently differ by 10--30\,ms. Using it as a
domestic ceiling is unsound. Worse, at national scale it is the wrong functional form:
regressing clean minimum RTT on straight-line distance over 346 domestic destinations
yields $R^2 = 0.013$ with a slightly negative slope. With a median domestic path of
$\approx$470\,km (4.6\,ms of propagation) against 25--35\,ms of fixed access overhead,
distance does not drive domestic latency, and no multiplier over distance can describe it.

\subsection{R0: the measurement can support a latency rule}

\textbf{Claim.} The RTT data is a measurement rather than an artefact of the measuring
process. \textbf{Test.} Two checks requiring no external reference: does RTT ever
\emph{decrease} with hop distance, which propagation cannot produce; and does a fixed
router return a stable value under repetition? \textbf{Gate.} If either fails, no latency
rule may be evaluated on that data; the measurement is repaired rather than the thresholds
retuned.

Our initial route sweep failed both, having emitted one packet per hop limit under 120-way
concurrency: 22\% of adjacent hop pairs showed RTT falling by more than 20\,ms with
increasing distance, and a single fixed router spanned 4--853\,ms. Re-measurement at 6--8
threads with 20--30 packets per address yields 1--2\,ms of interquartile dispersion and
values a median of 30--41\,ms lower. All latency figures below derive from the
re-measurement. \textbf{Routers are measured in-path}, with packets addressed to the
original destination under a constrained hop limit, never by addressing the router
directly: routers deprioritise traffic addressed to themselves, and the route \emph{to} a
router need not coincide with the route \emph{through} it.

\subsection{R1: the vantage's access chain}

\textbf{Claim.} All traces from vantage $v$ share an initial router sequence $C(v)$ until
routes diverge. \textbf{Test.} The dominant address at each early hop position.
Parameter-free, since the divergence point is measured.

At our vantage, three routers appear in $\geq$98.8\% of traces, diverging at hop~6, and
cost $B_{\mathrm{access}} = 20$\,ms. This is our own connection and is subtracted before
any latency is attributed to a destination.

\subsection{R2: an address is not where it is registered}

\textbf{Claim.} If address $X$ registered in country $C$ were physically in $C$, no
Pakistani vantage could observe it faster than $2\,d(v,C)/c_f$. \textbf{Test.} Clean
in-path minimum RTT against that floor, with $d$ computed from per-address geolocation and
the vantage's own geolocated position. No routing factor is applied, so a falsification is
a physical impossibility rather than a judgement. \textbf{Gate.} Foreign claims only;
domestic distances are too short for physics to discriminate.

The minimum is the correct estimator here because queuing, rate limiting and slow ICMP
generation are strictly additive, so the minimum of a controlled burst is the closest
available approach to propagation cost.

Of the foreign-registered addresses observed as hops, \textbf{79 are falsified}, 341 are
upheld, and 78 are undecidable because their country's floor lies inside the domestic
range. The falsified set includes US Department of Defense addresses responding in
$\approx$3\,ms, together with Cogent, Cloudflare, AT\&T and T-Mobile ranges: Pakistani
routers numbered from foreign address space. Absent this step these hops would be counted
as departures and would fabricate thousands of spurious detours.

\textbf{R2 falsifies a claim; it does not establish a location.} ``Not in Virginia'' is
not ``in Pakistan'', which is what R4$'$ decides.

\subsection{R3: the vantage can observe paths, under this protocol}

\textbf{Claim.} Traces from $v$ contain public addresses. \textbf{Gate.} Established per
protocol, never universally: a probe is usable \emph{under a protocol}, and in earlier
data the same device yielded no public hop in 100\% of TCP traces and 11\% of ICMP traces.
Here, 3 of 43{,}765 traces contained no public router.

\subsection{R4$'$: the vantage's domestic normal is its own}

\textbf{Claim.} A destination whose clean RTT lies far outside this vantage's own domestic
distribution did not traverse a domestic path.

\textbf{Test.} Construct the reference set from destinations whose observed path is
entirely domestic; measure each cleanly and retain the minimum; flag any value beyond the
Tukey far fence,
\begin{equation}
\tau \;=\; Q_3 + 3\,(Q_3 - Q_1).
\label{eq:fence}
\end{equation}

Equation~\eqref{eq:fence} is used because it is a standard, scale-free outlier definition
that no one selects to fit a dataset. A slow vantage obtains a high baseline \emph{and} a
high fence, so \textbf{the procedure transfers between vantages without retuning}, which a
millisecond threshold cannot.

\textbf{Gates.} At least $\approx$100 reference destinations, or the distribution's shape
is noise; and if more than half the reference set is itself detouring, the baseline
\emph{is} the detour and the rule is rejected for that vantage.

\textbf{Scope.} The reference set is selected topologically, so R4$'$ can only recover
detours the topological test missed. It is a second filter beneath the first, not a
replacement.

\begin{table}[t]
\centering
\caption{R4$'$ evaluated at our vantage. The bound is an output of the procedure, not a
constant of the method.}
\label{tab:r4}
\begin{tabular}{lr}
\toprule
Reference destinations              & 347 \\
Domestic median                     & 27\,ms \\
$Q_1$ / $Q_3$                       & 24 / 38\,ms \\
\textbf{Fence $\tau$}               & \textbf{80\,ms} \\
Domestic above $\tau$               & 1.4\% \\
Known detours below $\tau$          & 0 of 10 \\
\bottomrule
\end{tabular}
\end{table}

The two distributions separate empirically: the highest domestic value below the fence is
69\,ms and the lowest tromboning destination is 90\,ms.

\subsection{R5: the trace left the country and returned}

Evaluated only after R0--R4$'$ are settled for that vantage under that protocol. Anything
attributable to an earlier rule is subtracted first, and a detour is the residue:

\begin{enumerate}\itemsep2pt
\item subtract $C(v)$; the access chain cannot be a detour;
\item apply R2; a foreign registry is not evidence, an impossible RTT is;
\item apply R4$'$ to whatever R2 falsified, since ``not in $C$'' is not ``domestic'';
\item any surviving foreign hop constitutes a departure.
\end{enumerate}

A hop on a \emph{foreign internet exchange fabric} counts as a departure: the address
belongs to the exchange rather than to either peer, but the packet is physically abroad.

\textbf{Result: 930 of 34{,}191 traces (2.72\%), across 156 blocks in 49 networks.} This
is a \textbf{floor}. With the latency arm confined to R4$'$'s scope, a detour is visible
only when a foreign router responds; detours through silent routers or inside MPLS tunnels
are recorded as domestic.

\paragraph{Blocks, not traces, are the unit of evidence.} Tromboning is a property of the
destination prefix: 146 of 156 affected blocks are all-or-nothing, every selected host
detouring or none, and 143 of 156 exit via a single common address. Eight traces into one
block are eight observations of one routing decision. Confidence intervals must therefore
be computed over blocks and clustered by network.

\subsection{R6: the verdict predicts what it was not given}

\textbf{Claim.} If R5 detects a real phenomenon, its verdicts should separate on a
quantity the topological test never observed.

\textbf{Test.} Compare clean minimum RTT between destinations R5 labels domestic and those
it labels tromboning. Destinations labelled domestic have a median of \textbf{27\,ms};
those labelled tromboning, \textbf{143\,ms}. The topological classifier had no access to
latency, so the 116\,ms separation is independent corroboration rather than a restatement.

\textbf{Limit.} The tromboning arm of this comparison holds 10 destinations. It is
corroboration, not a significance test, and warrants repetition on a larger deliberate
sample.

\section{Limitations}
\label{sec:limits}

\textbf{Counts are floors with unequal tightness.} Adaptive sampling caps at 64 addresses
per block, so no block reports more than 64 regardless of occupancy. Independently,
re-probing recorded-dead addresses returns $\approx$35\% alive on one large network
against 22--25\% elsewhere. Densities are therefore comparable within a network over time,
not between networks, and no single national correction is admissible.

\textbf{Liveness is vantage-relative.} In a controlled retest, $\approx$9\% of addresses
confirmed alive from one Pakistani network did not respond from another, at every
concurrency tested. Reported counts are a union over two vantages covering
\emph{disjoint} address sets, not a two-vantage measurement of the same addresses.

\textbf{Route collection is effectively single-vantage.} 43{,}671 of 43{,}765 traces
egress one network. Reach rates are properties of that vantage's view, not of the
destination network alone.

\textbf{Proximate foreign destinations are undetectable by latency.} Domestic paths run at
a 27\,ms median while the physical floors to the nearest Gulf hubs are 4--7\,ms. A Gulf
detour and a slow domestic path are indistinguishable by RTT at any threshold; only
topology can recover them.

\textbf{Selection favours dense blocks.} A block requires live hosts before it can be
traced, and the retention gate then prefers responsive paths. A control comparison against
sparsely populated blocks has not been run, so the magnitude of this bias is unknown
rather than small.

\textbf{Remaining limits.} Traceroute observes the forward path only; return paths may
differ and are invisible. MPLS tunnels conceal interior hops. The study is IPv4 only,
covers a single country over three days, and no operator has confirmed any result.

%
%  Assumes the preamble already loads: booktabs, amsmath, amssymb.
%  Uses no algorithm package: the pseudocode sits in a plain figure float.
%  Cites: klick2016, bozkurt:dissecting:2018, ripestat  (all in references.bib)
%
%  Every figure below is produced by the pipeline in
%  experiments/15_national_block_routing_census/.
% =====================================================================

\section{Measurement Method}
\label{sec:method}

\subsection{Universe}
\label{sec:universe}

``Pakistan's address space'' has two definitions that disagree, and neither contains the
other. The \emph{registry view} is the space recorded against Pakistan at the RIR. The
\emph{routing view} is the space a Pakistani AS originates in BGP. Registration is
paperwork; origination is a deliberate operational act, and BGP does not require the two
to agree.

We take the union. Of the resulting space, \textbf{252{,}672 addresses (4.4\%) are
originated by a Pakistani AS but registered elsewhere}, typically leased space announced
as a more-specific inside a covering announcement the registrant still originates. In the
other direction, \textbf{1{,}209 /24 blocks} are registered to Pakistan and originated by
nobody. A study defining its universe by registration alone discards the first set; one
defining it by origination alone discards the second.

\begin{table}[t]
\centering
\caption{Measurement universe.}
\label{tab:universe}
\begin{tabular}{lr}
\toprule
Addresses (registry $\cup$ routing) & 5{,}774{,}336 \\
Announced or registered prefixes    & 932 \\
/24-equivalent blocks               & 22{,}556 \\
Networks with announced space       & 352 \\
\bottomrule
\end{tabular}
\end{table}

Prefixes are subdivided into /24-equivalent \emph{blocks} so that sampling density is
uniform: without this a /16 and a /24 would receive equal attention despite differing by a
factor of 256. The mapping is one-to-many, not one-to-one: 686 of the 932 prefixes are
larger than a /24. A per-block count is therefore not a per-prefix count.

Origination is queried through RIPEstat~\cite{ripestat}, which is fed by RIS. We note that
origination is \emph{not} a boolean: a prefix is originated \emph{as seen by $n$ of the
324 IPv4 RIS peers}, and RIPEstat endpoints that report a single visibility-filtered
origin can disagree with the full origin history for the same prefix at the same instant.
Attribution is therefore taken from origin history rather than from a filtered view.

\subsection{Adaptive liveness sampling}
\label{sec:sampling}

Exhaustive probing of 5.77M addresses is neither necessary nor courteous. We follow the
prefix-as-unit and density formulation of TASS~\cite{klick2016}, but \emph{invert its
rule}: TASS discards low-density prefixes to shrink a full-space scan, whereas we retain
every prefix and reduce effort within each one.

A \emph{check} tests one address by trying, in order, ICMP echo, TCP connect to port~80,
and TCP connect to port~443, stopping at the first response. A refused TCP connection
counts as alive, since something was present to refuse it. One check may thus emit up to
three packets. A \emph{draw} is a batch of $D=8$ addresses.

\begin{figure}[t]
\small
\hrule\vspace{4pt}
\begin{tabbing}
xx\=xx\=xx\=xx\kill
\textbf{for each} block $B$: \\
\> $S \leftarrow$ addresses of $B$ already probed by any prior scan \\
\> $L \leftarrow$ live hosts already known in $B$ \\
\> $P \leftarrow \mathrm{shuffle}\!\left(\mathrm{hosts}(B) \setminus S\right)$
   \quad \textit{(without replacement)} \\
\> $r \leftarrow 0$ \\
\> \textbf{while} $P \neq \emptyset$ \textbf{and} $|S| < M$ \textbf{and} $L < K$: \\
\> \> $b \leftarrow$ next $D$ addresses of $P$; \quad $r \leftarrow r+1$ \\
\> \> \textbf{for} $a \in b$: \\
\> \> \> $S \leftarrow S \cup \{a\}$;\quad
        \textbf{if} $\mathrm{check}(a)$ \textbf{then} $L \leftarrow L+1$ \\
\> \> \> record $(a, \mathrm{result})$ \quad \textit{(negative results retained)} \\
\> \> \textbf{if} $L = 0$ \textbf{and} $r \geq 2$ \textbf{then break} \\
\> emit $L$ for $B$
\end{tabbing}
\vspace{2pt}\hrule
\caption{Adaptive per-block sampling. $D=8$, $K=8$, $M=64$.}
\label{alg:sample}
\end{figure}

Sampling is \textbf{without replacement}: $S$ accumulates across draws and across scans,
so no address is ever probed twice. Draws accumulate rather than replace, and negative
results are retained because the miss-rate estimates of \S\ref{sec:limits} depend on them.
The target is re-tested per draw rather than per address, so a block may terminate above
$K$; 17.2\% of live blocks did.

\paragraph{Choice of $K$.} $K=8$ is measured rather than assumed, and the criterion is
\emph{split-block detection}, not redundancy. Over 1{,}758 (block, vantage) pairs that
received a full $K=8$ in a prior round, \textbf{20.0\% of blocks were split}: some probed
addresses detoured abroad while others remained domestic, so ``the route to a block'' is
not always well defined. On the 79.4\% of pairs that were unambiguous, a randomly chosen
$K=2$ reproduces the $K=8$ verdict 98.6\% of the time. A small $K$ therefore recovers the
\emph{verdict} while destroying the \emph{picture}: with two samples there is nothing to
compare against, and a split block is silently recorded as whichever pair was drawn. A
20\% split rate is too large to sample away.

\paragraph{Cost of the stopping rule.} The two-empty-draw rule is the design's principal
source of missed occupancy, and its cost is calculable. Sampling without replacement from
$N=254$ usable addresses, the probability that all $d=16$ draws miss a block containing
$\ell$ live hosts is hypergeometric:
\begin{equation}
\Pr[\text{all } d \text{ dead}] \;=\; \prod_{i=0}^{d-1} \frac{N-\ell-i}{N-i}.
\label{eq:hyper}
\end{equation}

\begin{table}[t]
\centering
\caption{Probability that the two-empty-draw rule abandons a populated block,
from Eq.~\eqref{eq:hyper}.}
\label{tab:hyper}
\begin{tabular}{rrr}
\toprule
Live hosts $\ell$ & Density & $\Pr[\text{abandoned}]$ \\
\midrule
1  & 0.4\%  & 93.7\% \\
2  & 0.8\%  & 87.8\% \\
5  & 2.0\%  & 72.0\% \\
8  & 3.1\%  & 59.0\% \\
13 & 5.1\%  & 42.0\% \\
26 & 10.2\% & 16.8\% \\
51 & 20.1\% & \phantom{0}2.4\% \\
\bottomrule
\end{tabular}
\end{table}

A block holding eight live hosts is abandoned 59\% of the time. \textbf{Block-level
coverage figures are therefore floors, and substantially so}: blocks we record as empty
are not reliably empty. This does not bias the routing results of \S\ref{sec:rules},
which condition on blocks that were successfully traced, but it forbids any claim about
the extent of unused address space.

\begin{table}[t]
\centering
\caption{Where sampling effort went. Effort follows signal: a populated block
receives roughly four times the attention of an empty one.}
\label{tab:effort}
\begin{tabular}{lrrr}
\toprule
 & Blocks & Checks & Mean/block \\
\midrule
Recorded empty & 16{,}669 & 364{,}972 & 21.9 \\
Contained life & \phantom{0}5{,}887 & 479{,}130 & 81.4 \\
\midrule
Total          & 22{,}556 & 844{,}102 & 37.4 \\
\bottomrule
\end{tabular}
\end{table}

\textbf{844{,}102 addresses were checked, 14.6\% of the universe, yielding 43{,}737 live
hosts.} The answer rate is 5.2\% \emph{of what was tested}; the remaining 4{,}930{,}234
addresses were never probed, and the count must not be read against the full space.

\subsection{Route collection and trace selection}
\label{sec:routes}

Every live host was traced with ICMP, incrementing the hop limit and terminating after
five consecutive silent hops. Of \textbf{43{,}765} traces attempted, \textbf{36{,}214
(82.7\%)} reached their target.

A trace is retained only if it \emph{reached} its target, so that the endpoint is known,
and at least five hops responded, so that the path has discernible shape. At most eight
traces are kept per block, ranked by fewest timeouts. This yields \textbf{34{,}191 traces
across 5{,}508 blocks in 217 networks}, comprising \textbf{217{,}255 responding hops},
each resolved to a country and an operator.

The five-hop gate is a judgement rather than a derived quantity, and it does not fall
evenly: it binds almost entirely on one large network whose paths are of ordinary length
but answer at half the usual rate (median 5 responding hops against 8 elsewhere, 48\%
silent slots against 31\%).

\section{Detection Rules}
\label{sec:rules}

A rule here is \textbf{a procedure, not a threshold}. Each is stated as a claim that could
be false, tested against a vantage point's own data, and applied only where it survives. A
rule proven at one vantage says nothing about another until retested there. Consequently
every numeric bound below is an \emph{output} of the procedure at our vantage, not a
constant of the method.

\paragraph{The single constant.} The speed of light in fibre, $c_f = 204{,}000$\,km/s,
taken from Bozkurt et al.~\cite{bozkurt:dissecting:2018}. It is a physical property, and
it appears in exactly one rule. A detector cannot derive physics from its own traffic.

We emphasise what is \emph{not} used. Bozkurt et al. also give a widely quoted rule of
thumb, that internet latency is approximately line-of-sight distance $\times\,2.1$ divided
by $c_f$. That is a \emph{typical-case estimator} and not an upper bound; the same work
reports that only 11\% of real fibre links fall within 25\% of the latency their length
predicts, and that co-located server pairs frequently differ by 10--30\,ms. Using it as a
domestic ceiling is unsound. Worse, at national scale it is the wrong functional form:
regressing clean minimum RTT on straight-line distance over 346 domestic destinations
yields $R^2 = 0.013$ with a slightly negative slope. With a median domestic path of
$\approx$470\,km (4.6\,ms of propagation) against 25--35\,ms of fixed access overhead,
distance does not drive domestic latency, and no multiplier over distance can describe it.

\subsection{R0: the measurement can support a latency rule}

\textbf{Claim.} The RTT data is a measurement rather than an artefact of the measuring
process. \textbf{Test.} Two checks requiring no external reference: does RTT ever
\emph{decrease} with hop distance, which propagation cannot produce; and does a fixed
router return a stable value under repetition? \textbf{Gate.} If either fails, no latency
rule may be evaluated on that data; the measurement is repaired rather than the thresholds
retuned.

Our initial route sweep failed both, having emitted one packet per hop limit under 120-way
concurrency: 22\% of adjacent hop pairs showed RTT falling by more than 20\,ms with
increasing distance, and a single fixed router spanned 4--853\,ms. Re-measurement at 6--8
threads with 20--30 packets per address yields 1--2\,ms of interquartile dispersion and
values a median of 30--41\,ms lower. All latency figures below derive from the
re-measurement. \textbf{Routers are measured in-path}, with packets addressed to the
original destination under a constrained hop limit, never by addressing the router
directly: routers deprioritise traffic addressed to themselves, and the route \emph{to} a
router need not coincide with the route \emph{through} it.

\subsection{R1: the vantage's access chain}

\textbf{Claim.} All traces from vantage $v$ share an initial router sequence $C(v)$ until
routes diverge. \textbf{Test.} The dominant address at each early hop position.
Parameter-free, since the divergence point is measured.

At our vantage, three routers appear in $\geq$98.8\% of traces, diverging at hop~6, and
cost $B_{\mathrm{access}} = 20$\,ms. This is our own connection and is subtracted before
any latency is attributed to a destination.

\subsection{R2: an address is not where it is registered}

\textbf{Claim.} If address $X$ registered in country $C$ were physically in $C$, no
Pakistani vantage could observe it faster than $2\,d(v,C)/c_f$. \textbf{Test.} Clean
in-path minimum RTT against that floor, with $d$ computed from per-address geolocation and
the vantage's own geolocated position. No routing factor is applied, so a falsification is
a physical impossibility rather than a judgement. \textbf{Gate.} Foreign claims only;
domestic distances are too short for physics to discriminate.

The minimum is the correct estimator here because queuing, rate limiting and slow ICMP
generation are strictly additive, so the minimum of a controlled burst is the closest
available approach to propagation cost.

Of the foreign-registered addresses observed as hops, \textbf{79 are falsified}, 341 are
upheld, and 78 are undecidable because their country's floor lies inside the domestic
range. The falsified set includes US Department of Defense addresses responding in
$\approx$3\,ms, together with Cogent, Cloudflare, AT\&T and T-Mobile ranges: Pakistani
routers numbered from foreign address space. Absent this step these hops would be counted
as departures and would fabricate thousands of spurious detours.

\textbf{R2 falsifies a claim; it does not establish a location.} ``Not in Virginia'' is
not ``in Pakistan'', which is what R4$'$ decides.

\subsection{R3: the vantage can observe paths, under this protocol}

\textbf{Claim.} Traces from $v$ contain public addresses. \textbf{Gate.} Established per
protocol, never universally: a probe is usable \emph{under a protocol}, and in earlier
data the same device yielded no public hop in 100\% of TCP traces and 11\% of ICMP traces.
Here, 3 of 43{,}765 traces contained no public router.

\subsection{R4$'$: the vantage's domestic normal is its own}

\textbf{Claim.} A destination whose clean RTT lies far outside this vantage's own domestic
distribution did not traverse a domestic path.

\textbf{Test.} Construct the reference set from destinations whose observed path is
entirely domestic; measure each cleanly and retain the minimum; flag any value beyond the
Tukey far fence,
\begin{equation}
\tau \;=\; Q_3 + 3\,(Q_3 - Q_1).
\label{eq:fence}
\end{equation}

Equation~\eqref{eq:fence} is used because it is a standard, scale-free outlier definition
that no one selects to fit a dataset. A slow vantage obtains a high baseline \emph{and} a
high fence, so \textbf{the procedure transfers between vantages without retuning}, which a
millisecond threshold cannot.

\textbf{Gates.} At least $\approx$100 reference destinations, or the distribution's shape
is noise; and if more than half the reference set is itself detouring, the baseline
\emph{is} the detour and the rule is rejected for that vantage.

\textbf{Scope.} The reference set is selected topologically, so R4$'$ can only recover
detours the topological test missed. It is a second filter beneath the first, not a
replacement.

\begin{table}[t]
\centering
\caption{R4$'$ evaluated at our vantage. The bound is an output of the procedure, not a
constant of the method.}
\label{tab:r4}
\begin{tabular}{lr}
\toprule
Reference destinations              & 347 \\
Domestic median                     & 27\,ms \\
$Q_1$ / $Q_3$                       & 24 / 38\,ms \\
\textbf{Fence $\tau$}               & \textbf{80\,ms} \\
Domestic above $\tau$               & 1.4\% \\
Known detours below $\tau$          & 0 of 10 \\
\bottomrule
\end{tabular}
\end{table}

The two distributions separate empirically: the highest domestic value below the fence is
69\,ms and the lowest tromboning destination is 90\,ms.

\subsection{R5: the trace left the country and returned}

Evaluated only after R0--R4$'$ are settled for that vantage under that protocol. Anything
attributable to an earlier rule is subtracted first, and a detour is the residue:

\begin{enumerate}\itemsep2pt
\item subtract $C(v)$; the access chain cannot be a detour;
\item apply R2; a foreign registry is not evidence, an impossible RTT is;
\item apply R4$'$ to whatever R2 falsified, since ``not in $C$'' is not ``domestic'';
\item any surviving foreign hop constitutes a departure.
\end{enumerate}

A hop on a \emph{foreign internet exchange fabric} counts as a departure: the address
belongs to the exchange rather than to either peer, but the packet is physically abroad.

\textbf{Result: 930 of 34{,}191 traces (2.72\%), across 156 blocks in 49 networks.} This
is a \textbf{floor}. With the latency arm confined to R4$'$'s scope, a detour is visible
only when a foreign router responds; detours through silent routers or inside MPLS tunnels
are recorded as domestic.

\paragraph{Blocks, not traces, are the unit of evidence.} Tromboning is a property of the
destination prefix: 146 of 156 affected blocks are all-or-nothing, every selected host
detouring or none, and 143 of 156 exit via a single common address. Eight traces into one
block are eight observations of one routing decision. Confidence intervals must therefore
be computed over blocks and clustered by network.

\subsection{R6: the verdict predicts what it was not given}

\textbf{Claim.} If R5 detects a real phenomenon, its verdicts should separate on a
quantity the topological test never observed.

\textbf{Test.} Compare clean minimum RTT between destinations R5 labels domestic and those
it labels tromboning. Destinations labelled domestic have a median of \textbf{27\,ms};
those labelled tromboning, \textbf{143\,ms}. The topological classifier had no access to
latency, so the 116\,ms separation is independent corroboration rather than a restatement.

\textbf{Limit.} The tromboning arm of this comparison holds 10 destinations. It is
corroboration, not a significance test, and warrants repetition on a larger deliberate
sample.

\section{Limitations}
\label{sec:limits}

\textbf{Counts are floors with unequal tightness.} Adaptive sampling caps at 64 addresses
per block, so no block reports more than 64 regardless of occupancy. Independently,
re-probing recorded-dead addresses returns $\approx$35\% alive on one large network
against 22--25\% elsewhere. Densities are therefore comparable within a network over time,
not between networks, and no single national correction is admissible.

\textbf{Liveness is vantage-relative.} In a controlled retest, $\approx$9\% of addresses
confirmed alive from one Pakistani network did not respond from another, at every
concurrency tested. Reported counts are a union over two vantages covering
\emph{disjoint} address sets, not a two-vantage measurement of the same addresses.

\textbf{Route collection is effectively single-vantage.} 43{,}671 of 43{,}765 traces
egress one network. Reach rates are properties of that vantage's view, not of the
destination network alone.

\textbf{Proximate foreign destinations are undetectable by latency.} Domestic paths run at
a 27\,ms median while the physical floors to the nearest Gulf hubs are 4--7\,ms. A Gulf
detour and a slow domestic path are indistinguishable by RTT at any threshold; only
topology can recover them.

\textbf{Selection favours dense blocks.} A block requires live hosts before it can be
traced, and the retention gate then prefers responsive paths. A control comparison against
sparsely populated blocks has not been run, so the magnitude of this bias is unknown
rather than small.

\textbf{Remaining limits.} Traceroute observes the forward path only; return paths may
differ and are invisible. MPLS tunnels conceal interior hops. The study is IPv4 only,
covers a single country over three days, and no operator has confirmed any result.

%
%  Assumes the preamble already loads: booktabs, amsmath, amssymb.
%  Uses no algorithm package: the pseudocode sits in a plain figure float.
%  Cites: klick2016, bozkurt:dissecting:2018, ripestat  (all in references.bib)
%
%  Every figure below is produced by the pipeline in
%  experiments/15_national_block_routing_census/.
% =====================================================================

\section{Measurement Method}
\label{sec:method}

\subsection{Universe}
\label{sec:universe}

``Pakistan's address space'' has two definitions that disagree, and neither contains the
other. The \emph{registry view} is the space recorded against Pakistan at the RIR. The
\emph{routing view} is the space a Pakistani AS originates in BGP. Registration is
paperwork; origination is a deliberate operational act, and BGP does not require the two
to agree.

We take the union. Of the resulting space, \textbf{252{,}672 addresses (4.4\%) are
originated by a Pakistani AS but registered elsewhere}, typically leased space announced
as a more-specific inside a covering announcement the registrant still originates. In the
other direction, \textbf{1{,}209 /24 blocks} are registered to Pakistan and originated by
nobody. A study defining its universe by registration alone discards the first set; one
defining it by origination alone discards the second.

\begin{table}[t]
\centering
\caption{Measurement universe.}
\label{tab:universe}
\begin{tabular}{lr}
\toprule
Addresses (registry $\cup$ routing) & 5{,}774{,}336 \\
Announced or registered prefixes    & 932 \\
/24-equivalent blocks               & 22{,}556 \\
Networks with announced space       & 352 \\
\bottomrule
\end{tabular}
\end{table}

Prefixes are subdivided into /24-equivalent \emph{blocks} so that sampling density is
uniform: without this a /16 and a /24 would receive equal attention despite differing by a
factor of 256. The mapping is one-to-many, not one-to-one: 686 of the 932 prefixes are
larger than a /24. A per-block count is therefore not a per-prefix count.

Origination is queried through RIPEstat~\cite{ripestat}, which is fed by RIS. We note that
origination is \emph{not} a boolean: a prefix is originated \emph{as seen by $n$ of the
324 IPv4 RIS peers}, and RIPEstat endpoints that report a single visibility-filtered
origin can disagree with the full origin history for the same prefix at the same instant.
Attribution is therefore taken from origin history rather than from a filtered view.

\subsection{Adaptive liveness sampling}
\label{sec:sampling}

Exhaustive probing of 5.77M addresses is neither necessary nor courteous. We follow the
prefix-as-unit and density formulation of TASS~\cite{klick2016}, but \emph{invert its
rule}: TASS discards low-density prefixes to shrink a full-space scan, whereas we retain
every prefix and reduce effort within each one.

A \emph{check} tests one address by trying, in order, ICMP echo, TCP connect to port~80,
and TCP connect to port~443, stopping at the first response. A refused TCP connection
counts as alive, since something was present to refuse it. One check may thus emit up to
three packets. A \emph{draw} is a batch of $D=8$ addresses.

\begin{figure}[t]
\small
\hrule\vspace{4pt}
\begin{tabbing}
xx\=xx\=xx\=xx\kill
\textbf{for each} block $B$: \\
\> $S \leftarrow$ addresses of $B$ already probed by any prior scan \\
\> $L \leftarrow$ live hosts already known in $B$ \\
\> $P \leftarrow \mathrm{shuffle}\!\left(\mathrm{hosts}(B) \setminus S\right)$
   \quad \textit{(without replacement)} \\
\> $r \leftarrow 0$ \\
\> \textbf{while} $P \neq \emptyset$ \textbf{and} $|S| < M$ \textbf{and} $L < K$: \\
\> \> $b \leftarrow$ next $D$ addresses of $P$; \quad $r \leftarrow r+1$ \\
\> \> \textbf{for} $a \in b$: \\
\> \> \> $S \leftarrow S \cup \{a\}$;\quad
        \textbf{if} $\mathrm{check}(a)$ \textbf{then} $L \leftarrow L+1$ \\
\> \> \> record $(a, \mathrm{result})$ \quad \textit{(negative results retained)} \\
\> \> \textbf{if} $L = 0$ \textbf{and} $r \geq 2$ \textbf{then break} \\
\> emit $L$ for $B$
\end{tabbing}
\vspace{2pt}\hrule
\caption{Adaptive per-block sampling. $D=8$, $K=8$, $M=64$.}
\label{alg:sample}
\end{figure}

Sampling is \textbf{without replacement}: $S$ accumulates across draws and across scans,
so no address is ever probed twice. Draws accumulate rather than replace, and negative
results are retained because the miss-rate estimates of \S\ref{sec:limits} depend on them.
The target is re-tested per draw rather than per address, so a block may terminate above
$K$; 17.2\% of live blocks did.

\paragraph{Choice of $K$.} $K=8$ is measured rather than assumed, and the criterion is
\emph{split-block detection}, not redundancy. Over 1{,}758 (block, vantage) pairs that
received a full $K=8$ in a prior round, \textbf{20.0\% of blocks were split}: some probed
addresses detoured abroad while others remained domestic, so ``the route to a block'' is
not always well defined. On the 79.4\% of pairs that were unambiguous, a randomly chosen
$K=2$ reproduces the $K=8$ verdict 98.6\% of the time. A small $K$ therefore recovers the
\emph{verdict} while destroying the \emph{picture}: with two samples there is nothing to
compare against, and a split block is silently recorded as whichever pair was drawn. A
20\% split rate is too large to sample away.

\paragraph{Cost of the stopping rule.} The two-empty-draw rule is the design's principal
source of missed occupancy, and its cost is calculable. Sampling without replacement from
$N=254$ usable addresses, the probability that all $d=16$ draws miss a block containing
$\ell$ live hosts is hypergeometric:
\begin{equation}
\Pr[\text{all } d \text{ dead}] \;=\; \prod_{i=0}^{d-1} \frac{N-\ell-i}{N-i}.
\label{eq:hyper}
\end{equation}

\begin{table}[t]
\centering
\caption{Probability that the two-empty-draw rule abandons a populated block,
from Eq.~\eqref{eq:hyper}.}
\label{tab:hyper}
\begin{tabular}{rrr}
\toprule
Live hosts $\ell$ & Density & $\Pr[\text{abandoned}]$ \\
\midrule
1  & 0.4\%  & 93.7\% \\
2  & 0.8\%  & 87.8\% \\
5  & 2.0\%  & 72.0\% \\
8  & 3.1\%  & 59.0\% \\
13 & 5.1\%  & 42.0\% \\
26 & 10.2\% & 16.8\% \\
51 & 20.1\% & \phantom{0}2.4\% \\
\bottomrule
\end{tabular}
\end{table}

A block holding eight live hosts is abandoned 59\% of the time. \textbf{Block-level
coverage figures are therefore floors, and substantially so}: blocks we record as empty
are not reliably empty. This does not bias the routing results of \S\ref{sec:rules},
which condition on blocks that were successfully traced, but it forbids any claim about
the extent of unused address space.

\begin{table}[t]
\centering
\caption{Where sampling effort went. Effort follows signal: a populated block
receives roughly four times the attention of an empty one.}
\label{tab:effort}
\begin{tabular}{lrrr}
\toprule
 & Blocks & Checks & Mean/block \\
\midrule
Recorded empty & 16{,}669 & 364{,}972 & 21.9 \\
Contained life & \phantom{0}5{,}887 & 479{,}130 & 81.4 \\
\midrule
Total          & 22{,}556 & 844{,}102 & 37.4 \\
\bottomrule
\end{tabular}
\end{table}

\textbf{844{,}102 addresses were checked, 14.6\% of the universe, yielding 43{,}737 live
hosts.} The answer rate is 5.2\% \emph{of what was tested}; the remaining 4{,}930{,}234
addresses were never probed, and the count must not be read against the full space.

\subsection{Route collection and trace selection}
\label{sec:routes}

Every live host was traced with ICMP, incrementing the hop limit and terminating after
five consecutive silent hops. Of \textbf{43{,}765} traces attempted, \textbf{36{,}214
(82.7\%)} reached their target.

A trace is retained only if it \emph{reached} its target, so that the endpoint is known,
and at least five hops responded, so that the path has discernible shape. At most eight
traces are kept per block, ranked by fewest timeouts. This yields \textbf{34{,}191 traces
across 5{,}508 blocks in 217 networks}, comprising \textbf{217{,}255 responding hops},
each resolved to a country and an operator.

The five-hop gate is a judgement rather than a derived quantity, and it does not fall
evenly: it binds almost entirely on one large network whose paths are of ordinary length
but answer at half the usual rate (median 5 responding hops against 8 elsewhere, 48\%
silent slots against 31\%).

\section{Detection Rules}
\label{sec:rules}

A rule here is \textbf{a procedure, not a threshold}. Each is stated as a claim that could
be false, tested against a vantage point's own data, and applied only where it survives. A
rule proven at one vantage says nothing about another until retested there. Consequently
every numeric bound below is an \emph{output} of the procedure at our vantage, not a
constant of the method.

\paragraph{The single constant.} The speed of light in fibre, $c_f = 204{,}000$\,km/s,
taken from Bozkurt et al.~\cite{bozkurt:dissecting:2018}. It is a physical property, and
it appears in exactly one rule. A detector cannot derive physics from its own traffic.

We emphasise what is \emph{not} used. Bozkurt et al. also give a widely quoted rule of
thumb, that internet latency is approximately line-of-sight distance $\times\,2.1$ divided
by $c_f$. That is a \emph{typical-case estimator} and not an upper bound; the same work
reports that only 11\% of real fibre links fall within 25\% of the latency their length
predicts, and that co-located server pairs frequently differ by 10--30\,ms. Using it as a
domestic ceiling is unsound. Worse, at national scale it is the wrong functional form:
regressing clean minimum RTT on straight-line distance over 346 domestic destinations
yields $R^2 = 0.013$ with a slightly negative slope. With a median domestic path of
$\approx$470\,km (4.6\,ms of propagation) against 25--35\,ms of fixed access overhead,
distance does not drive domestic latency, and no multiplier over distance can describe it.

\subsection{R0: the measurement can support a latency rule}

\textbf{Claim.} The RTT data is a measurement rather than an artefact of the measuring
process. \textbf{Test.} Two checks requiring no external reference: does RTT ever
\emph{decrease} with hop distance, which propagation cannot produce; and does a fixed
router return a stable value under repetition? \textbf{Gate.} If either fails, no latency
rule may be evaluated on that data; the measurement is repaired rather than the thresholds
retuned.

Our initial route sweep failed both, having emitted one packet per hop limit under 120-way
concurrency: 22\% of adjacent hop pairs showed RTT falling by more than 20\,ms with
increasing distance, and a single fixed router spanned 4--853\,ms. Re-measurement at 6--8
threads with 20--30 packets per address yields 1--2\,ms of interquartile dispersion and
values a median of 30--41\,ms lower. All latency figures below derive from the
re-measurement. \textbf{Routers are measured in-path}, with packets addressed to the
original destination under a constrained hop limit, never by addressing the router
directly: routers deprioritise traffic addressed to themselves, and the route \emph{to} a
router need not coincide with the route \emph{through} it.

\subsection{R1: the vantage's access chain}

\textbf{Claim.} All traces from vantage $v$ share an initial router sequence $C(v)$ until
routes diverge. \textbf{Test.} The dominant address at each early hop position.
Parameter-free, since the divergence point is measured.

At our vantage, three routers appear in $\geq$98.8\% of traces, diverging at hop~6, and
cost $B_{\mathrm{access}} = 20$\,ms. This is our own connection and is subtracted before
any latency is attributed to a destination.

\subsection{R2: an address is not where it is registered}

\textbf{Claim.} If address $X$ registered in country $C$ were physically in $C$, no
Pakistani vantage could observe it faster than $2\,d(v,C)/c_f$. \textbf{Test.} Clean
in-path minimum RTT against that floor, with $d$ computed from per-address geolocation and
the vantage's own geolocated position. No routing factor is applied, so a falsification is
a physical impossibility rather than a judgement. \textbf{Gate.} Foreign claims only;
domestic distances are too short for physics to discriminate.

The minimum is the correct estimator here because queuing, rate limiting and slow ICMP
generation are strictly additive, so the minimum of a controlled burst is the closest
available approach to propagation cost.

Of the foreign-registered addresses observed as hops, \textbf{79 are falsified}, 341 are
upheld, and 78 are undecidable because their country's floor lies inside the domestic
range. The falsified set includes US Department of Defense addresses responding in
$\approx$3\,ms, together with Cogent, Cloudflare, AT\&T and T-Mobile ranges: Pakistani
routers numbered from foreign address space. Absent this step these hops would be counted
as departures and would fabricate thousands of spurious detours.

\textbf{R2 falsifies a claim; it does not establish a location.} ``Not in Virginia'' is
not ``in Pakistan'', which is what R4$'$ decides.

\subsection{R3: the vantage can observe paths, under this protocol}

\textbf{Claim.} Traces from $v$ contain public addresses. \textbf{Gate.} Established per
protocol, never universally: a probe is usable \emph{under a protocol}, and in earlier
data the same device yielded no public hop in 100\% of TCP traces and 11\% of ICMP traces.
Here, 3 of 43{,}765 traces contained no public router.

\subsection{R4$'$: the vantage's domestic normal is its own}

\textbf{Claim.} A destination whose clean RTT lies far outside this vantage's own domestic
distribution did not traverse a domestic path.

\textbf{Test.} Construct the reference set from destinations whose observed path is
entirely domestic; measure each cleanly and retain the minimum; flag any value beyond the
Tukey far fence,
\begin{equation}
\tau \;=\; Q_3 + 3\,(Q_3 - Q_1).
\label{eq:fence}
\end{equation}

Equation~\eqref{eq:fence} is used because it is a standard, scale-free outlier definition
that no one selects to fit a dataset. A slow vantage obtains a high baseline \emph{and} a
high fence, so \textbf{the procedure transfers between vantages without retuning}, which a
millisecond threshold cannot.

\textbf{Gates.} At least $\approx$100 reference destinations, or the distribution's shape
is noise; and if more than half the reference set is itself detouring, the baseline
\emph{is} the detour and the rule is rejected for that vantage.

\textbf{Scope.} The reference set is selected topologically, so R4$'$ can only recover
detours the topological test missed. It is a second filter beneath the first, not a
replacement.

\begin{table}[t]
\centering
\caption{R4$'$ evaluated at our vantage. The bound is an output of the procedure, not a
constant of the method.}
\label{tab:r4}
\begin{tabular}{lr}
\toprule
Reference destinations              & 347 \\
Domestic median                     & 27\,ms \\
$Q_1$ / $Q_3$                       & 24 / 38\,ms \\
\textbf{Fence $\tau$}               & \textbf{80\,ms} \\
Domestic above $\tau$               & 1.4\% \\
Known detours below $\tau$          & 0 of 10 \\
\bottomrule
\end{tabular}
\end{table}

The two distributions separate empirically: the highest domestic value below the fence is
69\,ms and the lowest tromboning destination is 90\,ms.

\subsection{R5: the trace left the country and returned}

Evaluated only after R0--R4$'$ are settled for that vantage under that protocol. Anything
attributable to an earlier rule is subtracted first, and a detour is the residue:

\begin{enumerate}\itemsep2pt
\item subtract $C(v)$; the access chain cannot be a detour;
\item apply R2; a foreign registry is not evidence, an impossible RTT is;
\item apply R4$'$ to whatever R2 falsified, since ``not in $C$'' is not ``domestic'';
\item any surviving foreign hop constitutes a departure.
\end{enumerate}

A hop on a \emph{foreign internet exchange fabric} counts as a departure: the address
belongs to the exchange rather than to either peer, but the packet is physically abroad.

\textbf{Result: 930 of 34{,}191 traces (2.72\%), across 156 blocks in 49 networks.} This
is a \textbf{floor}. With the latency arm confined to R4$'$'s scope, a detour is visible
only when a foreign router responds; detours through silent routers or inside MPLS tunnels
are recorded as domestic.

\paragraph{Blocks, not traces, are the unit of evidence.} Tromboning is a property of the
destination prefix: 146 of 156 affected blocks are all-or-nothing, every selected host
detouring or none, and 143 of 156 exit via a single common address. Eight traces into one
block are eight observations of one routing decision. Confidence intervals must therefore
be computed over blocks and clustered by network.

\subsection{R6: the verdict predicts what it was not given}

\textbf{Claim.} If R5 detects a real phenomenon, its verdicts should separate on a
quantity the topological test never observed.

\textbf{Test.} Compare clean minimum RTT between destinations R5 labels domestic and those
it labels tromboning. Destinations labelled domestic have a median of \textbf{27\,ms};
those labelled tromboning, \textbf{143\,ms}. The topological classifier had no access to
latency, so the 116\,ms separation is independent corroboration rather than a restatement.

\textbf{Limit.} The tromboning arm of this comparison holds 10 destinations. It is
corroboration, not a significance test, and warrants repetition on a larger deliberate
sample.

\section{Limitations}
\label{sec:limits}

\textbf{Counts are floors with unequal tightness.} Adaptive sampling caps at 64 addresses
per block, so no block reports more than 64 regardless of occupancy. Independently,
re-probing recorded-dead addresses returns $\approx$35\% alive on one large network
against 22--25\% elsewhere. Densities are therefore comparable within a network over time,
not between networks, and no single national correction is admissible.

\textbf{Liveness is vantage-relative.} In a controlled retest, $\approx$9\% of addresses
confirmed alive from one Pakistani network did not respond from another, at every
concurrency tested. Reported counts are a union over two vantages covering
\emph{disjoint} address sets, not a two-vantage measurement of the same addresses.

\textbf{Route collection is effectively single-vantage.} 43{,}671 of 43{,}765 traces
egress one network. Reach rates are properties of that vantage's view, not of the
destination network alone.

\textbf{Proximate foreign destinations are undetectable by latency.} Domestic paths run at
a 27\,ms median while the physical floors to the nearest Gulf hubs are 4--7\,ms. A Gulf
detour and a slow domestic path are indistinguishable by RTT at any threshold; only
topology can recover them.

\textbf{Selection favours dense blocks.} A block requires live hosts before it can be
traced, and the retention gate then prefers responsive paths. A control comparison against
sparsely populated blocks has not been run, so the magnitude of this bias is unknown
rather than small.

\textbf{Remaining limits.} Traceroute observes the forward path only; return paths may
differ and are invisible. MPLS tunnels conceal interior hops. The study is IPv4 only,
covers a single country over three days, and no operator has confirmed any result.
